\documentclass[prb,preprint]{revtex4-1}
% AJP uses the same style as Phys. Rev. B (prb).

\usepackage{amsmath}
\usepackage{amsfonts}

\DeclareMathOperator\born{born}

\begin{document}

\title{Deriving quantum mechanics from probabilistic classical mechanics}

% Anonymous for review submission:
%\author{Warren MacEvoy}
%\affiliation{Department of Electrical and Computer Engineering, Colorado Mesa University, Grand Junction, CO 81501}

\date{\today}

\begin{abstract}
We show that non-relativistic quantum mechanics can be derived from classical
Hamiltonian dynamics with probabilistic initial data by lifting the equations of motion
to a Liouville equation, Fourier transforming in momentum, and applying a centered-difference
approximation.  The resulting equation separates into forward and backward
Schr\"odinger equations.  The only approximation is $O(\hbar^2)$ relative, rendering it
indistinguishable from exact.  The Born rule emerges as the instantaneous marginal
over momentum.  We discuss implications for the uncertainty principle, entanglement,
and the extension to general relativity via the equivalence principle.
\end{abstract}

\maketitle


\section{Introduction}

The goal of these notes is to show non-relativistic quantum mechanics can be closely
approximated by classical dynamics with probabilistic initial data, along with an interesting
transformation of coordinates.  The approximation may be so close as to be indistinguishable.
We give the notes in terms of one spatial dimension, but it is trivial to extend to higher
dimensions and multiple particles.


\section{Quantum perspective}

Consider a tiny one-dimensional non-relativistic particle of mass $M$ in a tiny potential $V(x)$.
Here the accepted model is the Schr\"odinger equation, where $|\Psi(x,t)|^2$ gives the probability
of observing the particle at location $x$ at time $t$.  The wave function $\Psi(x,t)$
evolves according to the partial differential equation
\begin{equation}\label{eq:schrodinger}
i\hbar \Psi_t = -\frac{\hbar^2}{2M}\Psi_{xx} + V(x)\Psi \,,
\end{equation}
where $\hbar \approx 1.05 \times 10^{-34}$~joule-seconds.


\section{Classical perspective}

Consider a regular one-dimensional non-relativistic particle of mass $M$ in a regular potential $V(x)$.
The position of the particle $x(t)$ and momentum of the particle $p(t)$ evolve according
to the ordinary differential equation
\begin{eqnarray}
M \frac{dx}{dt} & = & p \,, \label{eq:hamilton-x}\\
\frac{dp}{dt} & = & -\frac{\partial V}{\partial x} \,. \label{eq:hamilton-p}
\end{eqnarray}


\section{Classical from quantum perspective}

There should be a way to see classical mechanics arrive as an approximation of quantum
mechanics.  This should not be surprising, and is included here for completeness.\cite{landau1977}

Approximate
\begin{equation}
\Psi \approx A \exp\left(\frac{iS}{\hbar}\right) \,,
\end{equation}
and, keeping the highest order terms:
\begin{equation}
-S_t = \frac{1}{2M}\left(S_x\right)^2 + V \,.
\end{equation}
Differentiate with respect to $x$:
\begin{equation}
-S_{xt} = \frac{1}{M}S_x S_{xx} + \frac{\partial V}{\partial x} \,.
\end{equation}
Now define $P = S_x$
\begin{equation}
-P_t = \frac{1}{M}P P_x + \frac{\partial V}{\partial x} \,.
\end{equation}
This is a first order hyperbolic differential equation and is solvable by the method of
characteristics.  Using $p(t) = P(x(t),t)$:
\begin{equation}
\frac{dp}{dt} = \frac{\partial P}{\partial x}\frac{dx}{dt} + \frac{\partial P}{\partial t}
= \frac{\partial P}{\partial x}\frac{dx}{dt} - \frac{1}{M}P P_x - \frac{\partial V}{\partial x} \,.
\end{equation}
Choosing
\begin{eqnarray}
\frac{dx}{dt} &=& \frac{1}{M}p \,, \\
\frac{dp}{dt} &=& -\frac{\partial V}{\partial x}
\end{eqnarray}
reduces the equation.

So, to leading order, the Schr\"odinger equation carries phase-gradient
information along characteristics defined by the equations of motion for classical mechanics.
Note that characteristics will inevitably be singular.


\section{Three concepts to get quantum mechanics from classical mechanics}

What must be added to classical mechanics to get quantum mechanics?  We combine the following concepts.

\subsection{An ordinary differential equation can be lifted to a linear partial differential equation}

Think of an ordinary differential equation initial value problem as evolving a
partial differential equation which describes the evolution of the probability
distribution of the initial values.\cite{goldstein2002}

ODE:
\begin{eqnarray}
\frac{dx_1}{dt} &=& v_1 \,, \\
\frac{dx_2}{dt} &=& v_2 \,, \\
&\vdots& \nonumber \\
\frac{dx_n}{dt} &=& v_n \,.
\end{eqnarray}

PDE:
\begin{equation}\label{eq:liouville}
\frac{\partial w}{\partial t}
+ \frac{\partial (v_1 w)}{\partial x_1}
+ \frac{\partial (v_2 w)}{\partial x_2}
+ \cdots
+ \frac{\partial (v_n w)}{\partial x_n}
= 0 \,.
\end{equation}

Using the Dirac delta function for the initial value, the PDE recovers the ODE initial value
problem: $w(x,0) = w_0(x) = \delta(x-x_0)$.

\subsection{The Fourier transform}

\begin{equation}\label{eq:fourier-forward}
\hat{\Psi}(\hat{p}) = \int_{-\infty}^{+\infty} dp \, e^{-i\hat{p}p/\hbar}\,\Psi(p) \,,
\end{equation}
\begin{equation}\label{eq:fourier-inverse}
\Psi(p) = \int_{-\infty}^{+\infty} \frac{d\hat{p}}{2\pi \hbar} \, e^{+i\hat{p}p/\hbar}\,\hat{\Psi}(\hat{p}) \,,
\end{equation}
\begin{equation}
\partial_p \Psi \rightarrow i \frac{\hat{p}}{\hbar} \hat{\Psi} \,,
\end{equation}
\begin{equation}
p\Psi \rightarrow i \hbar \partial_{\hat{p}} \hat{\Psi} \,.
\end{equation}
\begin{equation}\label{eq:delta}
\delta(\hat{p})= \int_{-\infty}^{+\infty} \frac{dp}{2\pi \hbar} \, e^{i\hat{p}p/\hbar} \,.
\end{equation}
The use of $\hat{p} = \hbar \omega$ compared to the typical spectral parameter $\omega$ greatly
compresses the spectrum.

\subsection{An approximation of $V'$ by a centered difference}

\begin{equation}\label{eq:centered-diff}
\frac{\partial V(x)}{\partial x} \approx \frac{V(x+\varepsilon/2) - V(x-\varepsilon/2)}{\varepsilon}
+ O(\varepsilon^2)  \,.
\end{equation}
We actually approximate $\varepsilon V'$, so with first error term:
\begin{equation}\label{eq:centered-diff-error}
\varepsilon \frac{\partial V(x)}{\partial x} \approx V(x+\varepsilon/2) - V(x-\varepsilon/2)
+ \frac{\varepsilon^3}{24} \frac{\partial^3 V(x)}{\partial x^3} + O(\varepsilon^5)  \,.
\end{equation}


\section{Quantum from probabilistic classical mechanics}

We note the structure here is related to the Wigner-Moyal formulation,\cite{wigner1932,moyal1949,hillery1984}
arrived at here from the classical rather than quantum direction.

Using the ODE $\Rightarrow$ PDE concept for the classical Hamiltonian system, the ODE is
\begin{eqnarray}
\frac{dx}{dt} &=& \frac{1}{M}p \,, \\
\frac{dp}{dt} &=& -\frac{\partial V}{\partial x} \,.
\end{eqnarray}

This corresponds to the PDE describing the evolution of the probability density:
\begin{equation}
\frac{\partial w}{\partial t}
+ \frac{\partial}{\partial x}\left(\frac{1}{M}pw\right)
+ \frac{\partial}{\partial p}\left(-\frac{\partial V}{\partial x}w\right)
= 0 \,,
\end{equation}
or
\begin{equation}\label{eq:liouville-hamiltonian}
w_t + \frac{p}{M} w_x - \frac{\partial V}{\partial x} w_p = 0 \,.
\end{equation}

Now Fourier transform only with respect to the momentum $p$.  This results in
\begin{equation}\label{eq:fourier-liouville}
i\hbar \hat{w}_t - \frac{\hbar^2}{M} \hat{w}_{x\hat{p}} + \hat{p} \frac{\partial V}{\partial x} \hat{w} = 0 \,.
\end{equation}

We now change variables and use the only approximation.  Define
\begin{equation}
z^{\pm} = x \pm \frac{\hat{p}}{2} \,,
\end{equation}
so
\begin{eqnarray}
\frac{\partial}{\partial x} & = & \frac{\partial}{\partial z^+} + \frac{\partial}{\partial z^-}\,,\\
\frac{\partial}{\partial \hat{p}} &=& \frac{1}{2}\frac{\partial}{\partial z^+}
- \frac{1}{2}\frac{\partial}{\partial z^-}
\end{eqnarray}
and use the approximation~(\ref{eq:centered-diff}) with $\varepsilon = \hat{p}$:
\begin{equation}\label{eq:approx-V}
\frac{\partial V}{\partial x}
\approx \frac{V(x+\hat{p}/2) - V(x-\hat{p}/2)}{\hat{p}} \,.
\end{equation}

This transforms Eq.~(\ref{eq:fourier-liouville}) into
\begin{equation}\label{eq:w-hat-pde}
i\hbar \hat{w}_t - \frac{\hbar^2}{2M}
\left[\left(\frac{\partial}{\partial z^{+}}\right)^2 - \left(\frac{\partial}{\partial z^{-}}\right)^2\right]\hat{w}
+ [V(z^+) - V(z^-)]\hat{w} = 0 \,.
\end{equation}
While this does not yet look like Schr\"odinger's equation, this is a linear differential
equation, and we can look at separation of variables to construct a general solution
to Eq.~(\ref{eq:w-hat-pde}).
Suppose $\hat{w} = \sum \psi^+(z^+)\psi^-(z^-)$ where each of these terms satisfy the above; we have
\begin{equation}\label{eq:separated}
\mp i\hbar \psi_t^{\pm} = -\frac{\hbar^2}{2M}\psi^{\pm}_{zz} + V\psi^{\pm} \,.
\end{equation}
Here $z$ takes the respective roles of $z^+$ and $z^-$.  For real potential $V$, we can have
$\psi^+ = e^{i\lambda t} \psi_\lambda$ and $\psi^- = e^{-i\lambda t}\psi_\lambda^*$ where the $\psi_\lambda$ are eigenvectors
of the original Schr\"odinger equation~(\ref{eq:schrodinger}).  These form a complete basis for the
solution of PDE for $\hat{w}$:
\begin{equation}\label{eq:w-hat-expansion}
\hat{w}(t,z_+,z_-)=\sum_{\lambda,\lambda'} C_{\lambda,\lambda'} e^{i(\lambda-\lambda')t} \psi_{\lambda}(z_+) \psi^*_{\lambda'}(z_-) \,,
\end{equation}
where the $C_{\lambda,\lambda'}$ are determined by the initial distribution $w_0$.
Assuming $\psi_\lambda$ are normalized by $\int_{-\infty}^{+\infty} dz \, |\psi_\lambda(z)|^2 = 1$, we have:
\begin{eqnarray}
C_{\lambda,\lambda'} &=& \int_{-\infty}^{+\infty} dz^+ \int_{-\infty}^{+\infty} dz^- \, \psi_\lambda^*(z^+) \psi_{\lambda'}(z^-) \hat{w}_0(z^+,z^-) \,, \label{eq:C-z}\\
&=& \int_{-\infty}^{+\infty} dx \int_{-\infty}^{+\infty} d\hat{p} \, \psi_\lambda^*\!\left(x+\frac{\hat{p}}{2}\right) \psi_{\lambda'}\!\left(x-\frac{\hat{p}}{2}\right) \hat{w}_0(x,\hat{p}) \,. \label{eq:C-xp}
\end{eqnarray}
Note $C$ is Hermitian and $w$ is a probability distribution.  This implies $C=\sum_k c_k a_k a_k^\dagger$ for real eigenvalues $c_k$, unit eigenvectors $a_k$, and $\sum_k c_k = 1$.


\section{The Born projection}

The standard Born rule states that at time $t$, the probability density of observing position $x$ is $|\Psi(x,t)|^2$.  We recover this as the instantaneous marginal of $w$ over momentum, which we call the Born projection:
\begin{eqnarray}
\born w &=& \int_{-\infty}^{+\infty} dp \, w(x,p,t) \,, \label{eq:born-1}\\
&=& \int_{-\infty}^{+\infty} dp \int_{-\infty}^{+\infty} \frac{d\hat{p}}{2\pi \hbar} \, e^{i\hat{p}p/\hbar} \hat{w} \,, \label{eq:born-2}\\
&=& \int_{-\infty}^{+\infty} d\hat{p} \, \hat{w} \int_{-\infty}^{+\infty} \frac{dp}{2\pi \hbar} \, e^{i\hat{p}p/\hbar} \,, \label{eq:born-3}\\
&=& \int_{-\infty}^{+\infty} d\hat{p} \, \hat{w}\,\delta(\hat{p}) \,, \label{eq:born-4}\\
&=& \hat{w}|_{\hat{p}=0} \,, \label{eq:born-5}\\
&=& \sum_{\lambda,\lambda'} C_{\lambda,\lambda'} e^{i(\lambda-\lambda')t} \psi_\lambda(x)\psi^*_{\lambda'}(x) \,, \label{eq:born-6}\\
&=& \sum_k c_k \sum_{\lambda,\lambda'} (a_k)_{\lambda} (a_k)_{\lambda'}^* e^{i(\lambda-\lambda')t} \psi_\lambda(x)\psi^*_{\lambda'}(x) \,, \label{eq:born-7}\\
&=& \sum_k c_k \left| \sum_{\lambda} (a_k)_{\lambda} e^{i\lambda t} \psi_\lambda(x) \right|^2 \,. \label{eq:born-rule}
\end{eqnarray}

This is the standard Born rule of quantum mechanics, now derived from the classical probability dynamics.  Notice this does not affect the dynamics---it is merely the observable projection of the dynamics to a measurement.  The Born rule emerges without decoherence;\cite{zurek2003} measurement is simply the marginalization of a classical probability distribution over the unobserved variable.


\section{Demystifications}

\subsection{Heisenberg's uncertainty principle}

There is nothing wrong with a distribution that specifies the precise initial position and
momentum of the particle.  However, we can only observe the marginal over momentum, which makes
the view fuzzy.

Substitution of $w_0(x,p)=\delta(x-x_0)\delta(p-p_0)$ into the formula for $C_{\lambda,\lambda'}$
gives the cross-Wigner function as the coefficients of the expansion of $\hat{w}$ in the
eigenbasis of the Schr\"odinger equation:
\begin{equation}\label{eq:cross-wigner}
C_{\lambda,\lambda'}[\delta(x-x_0)\delta(p-p_0)] = \int_{-\infty}^{+\infty} d\hat{p} \, \psi_\lambda^*\!\left(x_0+\frac{\hat{p}}{2}\right) \psi_{\lambda'}\!\left(x_0-\frac{\hat{p}}{2}\right) e^{-i\hat{p}p_0/\hbar}\,.
\end{equation}

For a symmetric potential $V$, the eigenfunctions $\psi_\lambda$ are either symmetric or antisymmetric, $\psi_\lambda(-x)=\sigma_\lambda \psi_\lambda(x)$, where $\sigma_\lambda = \pm 1$.  We can use this to compute $C_{\lambda,\lambda'}$ for the $x_0=0$, $p_0=0$ case:
\begin{eqnarray}
  C_{\lambda,\lambda'}[\delta(x)\delta(p)] &=& \int_{-\infty}^{+\infty} d\hat{p} \, \psi_\lambda^*\!\left(+\frac{\hat{p}}{2}\right) \psi_{\lambda'}\!\left(-\frac{\hat{p}}{2}\right) \,, \\
  &=& \sigma_{\lambda'} \int_{-\infty}^{+\infty} d\hat{p} \, \psi_\lambda^*\!\left(\frac{\hat{p}}{2}\right) \psi_{\lambda'}\!\left(\frac{\hat{p}}{2}\right) \,, \\
  &=& 2 \sigma_{\lambda'} \int_{-\infty}^{+\infty} d u \, \psi_\lambda^*(u) \psi_{\lambda'}(u) \,,\\
  &=& 2 \sigma_\lambda \delta_{\lambda,\lambda'} \,.
\end{eqnarray}

The Born projection yields
\begin{equation}
  \born \delta(x) \delta (p) = 2 \sum_{\lambda} \sigma_{\lambda} |\psi_{\lambda}(x)|^2 \,.
\end{equation}
This is the expansion of $\delta(x)$ as a distribution---what we would expect for a particle at the bottom of the well with no momentum.

\subsection{Quantum entanglement}

The above derivation is for a single particle in a single space dimension, but it is trivial to extend to multiple particles in three spatial dimensions.  Entanglement amounts to a dependent choice of initial distribution, so $\hat{w}$ cannot be factored into independent components for each particle.  The underlying mechanism evolves independently; it is only the initial distribution and the momentum marginal that make it look like there is something mysterious.  It is a consequence of dependent distributions and the projection.

For two particles with single-particle eigenstates $\hat{\alpha}=\psi_{\lambda}$ and $\hat{\beta}=\psi_\mu$, the coefficient matrix $C$ extends to the two-particle eigenbasis $\{|\alpha\alpha\rangle, |\alpha\beta\rangle, |\beta\alpha\rangle, |\beta\beta\rangle\}$.  The four maximally entangled Bell states are pure states $C = a\,a^\dagger$ with
\begin{eqnarray}
a_{\Phi^{\pm}} &=& \tfrac{1}{\sqrt{2}}(|\alpha\alpha\rangle \pm |\beta\beta\rangle) \,,\\
a_{\Psi^{\pm}} &=& \tfrac{1}{\sqrt{2}}(|\alpha\beta\rangle \pm |\beta\alpha\rangle) \,.
\end{eqnarray}
None of these factor as $|s_1\rangle \otimes |s_2\rangle$; entanglement is precisely this non-factorizability of $a$.  The outer product $a\,a^\dagger$ generates off-diagonal coherence terms in $\hat{w}$ that distinguish entangled states from classical mixtures, which would have $C$ diagonal.

\subsection{Many worlds}

The many-worlds interpretation\cite{everett1957} posits a proliferation of branching universes to avoid the measurement problem.  In the present framework, there is no measurement problem to solve.  The underlying dynamics are deterministic classical Hamiltonian evolution of a probability distribution over phase space.  The apparent indeterminacy of quantum mechanics arises entirely from the Born projection---the momentum marginal---which discards information the observer cannot access.  This is ordinary classical probability: a single world with uncertain initial conditions, viewed through an isolating projection.  There is no need for many worlds when one world with indeterminacy suffices.

\subsection{General relativity}

The equivalence principle guarantees that spacetime is locally flat: in any sufficiently small region, a freely falling observer experiences ordinary non-relativistic mechanics.  In such a local frame, classical Hamiltonian dynamics applies, and the derivation above produces the Schr\"odinger equation with an $O(\hbar^2)$ relative error.  Gravitational effects enter as smooth potentials in the local frame, for which the centered-difference approximation introduces errors on the order of $\hbar^2 \approx 10^{-68}$ in SI units.  The Born rule remains the momentum marginal.  Non-relativistic quantum mechanics therefore emerges naturally in any locally flat patch of spacetime, inheriting its domain of validity from the local flatness of the geometry, without additional postulates.

\subsection{Quantum computation}

The present framework suggests that quantum computation may be more amenable to classical simulation than previously appreciated.  Because the dynamics are classical Hamiltonian evolution of a probability distribution over phase space, Monte Carlo sampling provides a natural simulation strategy.  For an $n$-particle system, the phase space is $\mathbb{R}^{6n}$, but Monte Carlo convergence is $O(1/\sqrt{N})$ independent of dimension, requiring $O(1/\varepsilon^2)$ samples for precision $\varepsilon$.  The implications of this observation will be addressed separately.


\section{Discussion}

\subsection{Reframing the error}

The standard reading of the Wigner-Moyal formulation\cite{wigner1932,moyal1949} begins with classical mechanics---the Liouville equation in phase space---and adds quantum corrections order by order in $\hbar^2$.  The first such correction,
\begin{equation}\label{eq:moyal-correction}
\frac{\hbar^2}{24}\frac{\partial^3 V}{\partial x^3}\frac{\partial^3 w}{\partial p^3} \,,
\end{equation}
is conventionally interpreted as the leading quantum effect absent from classical dynamics.

The derivation presented here inverts this narrative.  The classical Liouville equation, once Fourier-transformed in momentum and subjected to a centered-difference change of variables, \emph{already is} the Schr\"odinger equation.  Quantum mechanics is not the result of adding corrections to the classical system---it is what the classical probability dynamics looks like when viewed in $\hat{p}$ coordinates.

What the Moyal series calls the ``first quantum correction'' is therefore not a physical effect being added.  It is the residual error of the centered-difference approximation---an artifact of the coordinate transformation, not of the physics.  Crucially, this residual is an $O(\hbar^2)$ \emph{relative} correction to a quantity that is already $O(\hbar)$, giving an $O(\hbar^3)$ absolute error, on the order of $10^{-100}$ in SI units.  The ``quantum correction'' was never a correction \emph{toward} quantum mechanics.  It is a negligible coordinate artifact.

\subsection{The approximation is exact for quadratic potentials}

The centered-difference approximation is exact whenever $V'''(x) = 0$, which includes the free particle and the harmonic oscillator---potentials up to quadratic order.  For all other potentials, the error is $O(\hbar^2)$ relative, which is already academic.

\subsection{Non-smooth potentials}

For potentials where $V'''$ is large or singular, the centered difference effectively smooths $V$ on a scale of $\hbar^2$.  No measurement can resolve spatial features at this scale, so the smoothed potential is experimentally indistinguishable from the original.

\subsection{Implications}

The Moyal series has been known since 1949, but it has always been read left to right: classical mechanics, corrected toward quantum mechanics.  The present derivation reads it right to left, and from that direction, the corrections are demoted from fundamental physics to coordinate artifacts.  The mystery is not why quantum mechanics departs from classical mechanics.  It is why the universe presents its dynamics to us through a momentum marginal of a Fourier-transformed phase space, revealing only the position projection to observers.

Regarding Bell's theorem,\cite{bell1964,bell1987} note that the introduced hidden variables (classical momentum) are natural from the classical perspective, and the Fourier transform is the non-locality which allows this reformulation to co-exist with Bell's theorem.  The mathematics is indifferent to the interpretation---the formulations are effectively equivalent---but the interpretation is compelling from the simplicity of the underlying mechanism.

These notes stayed in a drawer for a long time, since I was not convinced they explained anything.  But:
\begin{itemize}
\item They provide a useful demystification of quantum mechanics, allowing the intuition of classical mechanics to provide insight.
\item They place the mystery specifically: why does the universe conspire to transform the momentum space, and choose to reveal only the position projection to observers?
\item They reframe the Moyal series: what was thought of as the first quantum correction to classical mechanics is actually a negligible relative error in a coordinate transformation that already contains the full quantum theory.
\item They suggest a global unknowable state dependence through the transform that survives observation.  Decoherence is real, but perhaps a side-effect of large interactions and not as central as postulated.\cite{zurek2003}
\end{itemize}

I hope others find these ideas interesting.


\begin{acknowledgments}
The author has no conflicts to disclose.
\end{acknowledgments}


\begin{thebibliography}{99}

\bibitem{landau1977} L. D. Landau and E. M. Lifshitz, \textit{Quantum Mechanics:
Non-Relativistic Theory}, 3rd edition (Pergamon Press, Oxford, 1977).

\bibitem{goldstein2002} Herbert Goldstein, Charles Poole, and John Safko,
\textit{Classical Mechanics}, 3rd edition (Addison-Wesley, San Francisco, 2002).

\bibitem{wigner1932} E. P. Wigner, ``On the quantum correction for thermodynamic
equilibrium,'' Phys. Rev. \textbf{40}, 749--759 (1932).

\bibitem{moyal1949} J. E. Moyal, ``Quantum mechanics as a statistical theory,''
Math. Proc. Cambridge Philos. Soc. \textbf{45}, 99--124 (1949).

\bibitem{hillery1984} M. Hillery, R. F. O'Connell, M. O. Scully, and E. P. Wigner,
``Distribution functions in physics: Fundamentals,''
Phys. Rep. \textbf{106} (3), 121--167 (1984).

\bibitem{zurek2003} W. H. Zurek, ``Decoherence, einselection, and the quantum
origins of the classical,'' Rev. Mod. Phys. \textbf{75} (3), 715--775 (2003).

\bibitem{everett1957} H. Everett III, ```Relative state' formulation of quantum mechanics,''
Rev. Mod. Phys. \textbf{29} (3), 454--462 (1957).

\bibitem{bell1964} J. S. Bell, ``On the Einstein Podolsky Rosen paradox,''
Physics \textbf{1} (3), 195--200 (1964).

\bibitem{bell1987} J. S. Bell, \textit{Speakable and Unspeakable in Quantum Mechanics}
(Cambridge University Press, Cambridge, 1987).

\end{thebibliography}

\end{document}
